% Part: incompleteness
% Chapter: theories-computability
% Section: first-incompleteness

\documentclass[../../../include/open-logic-section]{subfiles}

\begin{document}

\olfileid[pt]{inc}{tcp}{inc}

\olsection{$\Th{Q}$ não possui extensões completas, consistentes e
  !!^{axiomatizable}}

\begin{thm}
\ollabel{thm:first-incompleteness}
Não existe extensão completa, consistente e !!{axiomatizable} de
$\Th{Q}$.
\end{thm}

\begin{proof}
Já sabemos que não existe extensão consistente e decidível de
$\Th{Q}$. Mas, se $\Th{T}$ é completa e !!{axiomatized}, então é
decidível.
\end{proof}

\begin{explain}
Este teorema não está tão distante da formulação original, de 1931, do
Primeiro Teorema da Incompletude de G\"odel. Além da terminologia mais
moderna, as diferenças principais são estas: G\"odel usa
``$\omega$-consistente'' em vez de ``consistente''; e não podia dizer
``!!{axiomatizable}'' em toda a generalidade, pois a noção formal de
computabilidade ainda não estava estabelecida. (Os modelos formais de
computabilidade foram desenvolvidos ao longo da década seguinte,
inclusive por G\"odel, em grande parte para permitir a caracterização
dos tipos de teoria suscetíveis ao fenômeno de G\"odel.)

O teorema diz que não se pode ter tudo, a saber, completude,
consistência e !!{axiomatizability}. Se abrirmos mão de qualquer uma
delas, porém, podemos ter as outras duas: $\Th{Q}$ é consistente e
axiomatizada computavelmente, mas não completa; a teoria inconsistente
é completa e axiomatizada computavelmente (digamos, por
$\{ \eq/[0][0] \}$), mas não consistente; e o conjunto das sentenças
verdadeiras da aritmética é completo e consistente, mas não é
axiomatizado computavelmente.
\end{explain}
\end{document}
